% ============================================================================
% sec03.tex — §3 Thermal Foundations, Systems, Zeroth Law, Thermometry Scales,
%             Thermal Expansion & The 8 Named Scientists (2027 Rule 5.c.v)
% Target Page Budget: 8 pages | Figures: fig:03:scales, fig:03:bimetallic, fig:03:history
% ============================================================================
\binderchapter{3}{Thermal Foundations, Zeroth Law, Thermometry Scales, Thermal Expansion \& The 8 Named Scientists}{Foundations, Scales, Expansion \& 8 Scientists}{sec:03}
\scopebanner{REQUIRED}[2027 Rules 5.c.i, 5.c.v (All Tournament Levels --- Minimum 6 Written Questions Combined)]

\idxtopic{intensive vs.\ extensive properties}
\idxtopic{Zeroth Law of Thermodynamics}
\idxtopic{temperature scales!all 8 modern and historical scales}
\idxtopic{thermal expansion!linear, area, volume, bimetallic, stress}
\idxtopic{scientists (Rule 5.c.v)!Kelvin, Black, Joule, Carnot, Planck, Clausius, Boltzmann, Maxwell}

\begin{formulacard}[§3 Foundations, 8 Temperature Scales, Thermal Expansion \& 8 Scientists Quick Card]
\begin{ftab}
$\dfrac{T - T_1}{T_2 - T_1} = \dfrac{X - X_1}{X_2 - X_1}$ & $X$: linear thermometric property ($L, P_V, R$); $(X_1,T_1), (X_2,T_2)$: fixed points & \RULEf & Universal two-point linear thermometer calibration; gas thermometer $T = \SI{273.16}{\kelvin}(P/P_{\mathrm{tp}})$ \src{S13} \\
$\dfrac{T_C}{100} = \dfrac{T_F - 32}{180} = \dfrac{T_K - 273.15}{100}$ & $T_C$: Celsius; $T_F$: Fahrenheit; $T_K$: Kelvin; Rankine $T_R = 1.8\,T_K = T_F + 459.67$ & \RULEf & Primary four-scale conversions; interval rule $\Delta T_K = \Delta T_C = \tfrac{5}{9}\Delta T_F = \tfrac{5}{9}\Delta T_R$ \src{S08,S13} \\
$T_{\mathrm{R\acute{e}}} = 0.8\,T_C;\quad T_{\mathrm{De}} = 150 - 1.5\,T_C$ & $T_{\mathrm{R\acute{e}}}$: R\'eaumur ($0^\circ/80^\circ$); $T_{\mathrm{De}}$: Delisle (\textbf{inverted} $150^\circ/0^\circ$); $T_N = 0.33\,T_C$; $T_{\text{R\o}} = 0.525\,T_C + 7.5$ & \RULEf & Four historical competition scales (R\'eaumur, Delisle, Newton, R\o mer) \src{S25,S26} \\
$\Delta L = \alpha_L L_0 \Delta T;\quad \Delta A \approx 2\alpha_L A_0 \Delta T$ & $\alpha_L$: linear expansion coeff.\ $(\si{\per\kelvin})$; $L_0, A_0$: initial length $(\si{\meter})$, area $(\si{\square\meter})$ & \INCOMP & 1D rod/wire/hole/perimeter \& 2D isotropic plate/circular aperture expansion \src{S13,S14} \\
$\Delta V = \beta_V V_0 \Delta T \approx 3\alpha_L V_0 \Delta T$ & $\beta_V$: volume expansion coeff.\ $(\si{\per\kelvin})$; $\Delta V_{\mathrm{over}} = (\beta_{\mathrm{liq}} - 3\alpha_{\mathrm{ves}})V_0\Delta T$ & \INCOMP & 3D solid/cavity/liquid expansion \& apparent liquid overflow from expanding vessel \src{S13} \\
$\sigma = \dfrac{F}{A} = -E\alpha_L\Delta T;\quad R_{\mathrm{curv}} \approx \dfrac{d}{(\alpha_2-\alpha_1)\Delta T}$ & $E$: Young's modulus $(\si{\pascal})$; $d$: layer center spacing $(\si{\meter})$; $\alpha_2 > \alpha_1$ & \INCOMP & Clamped-rod thermal stress \& bimetallic strip radius of curvature \src{S13,S26} \\
\end{ftab}
\end{formulacard}

\section{Intensive vs.\ Extensive Properties, Equilibrium \& The Zeroth Law}
\label{sec:03.1}
\idxtopic{state functions vs.\ path functions}
\idxtopic{equilibrium!thermal, mechanical, phase, chemical vs.\ steady state}
\idxtopic{quasistatic vs.\ reversible processes}

\subsection{The Duplication Test for Intensive vs.\ Extensive Properties}
Consider combining two identical homogeneous systems, each in equilibrium state $(T, P, V, m, n, U, H, S, \rho, c_p)$, by removing a partition between them \src{S13,S16}:
\begin{itemize}
  \item \textbf{Extensive Properties (Scale with System Size, $Y \to 2Y$; Additive $Y_{\mathrm{tot}} = \sum Y_i$):} Mass $m$ $(\si{\kilogram})$, amount of substance $n$ $(\si{\mole})$, particle count $N$, volume $V$ $(\si{\cubic\meter})$, internal energy $U$ $(\si{\joule})$, enthalpy $H = U + PV$ $(\si{\joule})$, Helmholtz free energy $F = U - TS$ $(\si{\joule})$, Gibbs free energy $G = H - TS$ $(\si{\joule})$, entropy $S$ $(\si{\joule\per\kelvin})$, total heat capacity $C = mc = nC_m$ $(\si{\joule\per\kelvin})$, and parallel thermal conductance $G_{\mathrm{th}} = 1/R_{\mathrm{th}}$ $(\si{\watt\per\kelvin})$.
  \item \textbf{Intensive Properties (Independent of System Size, $y \to y$; Spatially Uniform at Equilibrium):} Temperature $T$ $(\si{\kelvin})$, pressure $P$ $(\si{\pascal})$, density $\rho = m/V$ $(\si{\kilogram\per\cubic\meter})$, specific volume $v = V/m = 1/\rho$ $(\si{\cubic\meter\per\kilogram})$, molar volume $V_m = V/n$ $(\si{\cubic\meter\per\mole})$, specific heat capacity $c = C/m$ $(\si{\joule\per\kilogram\per\kelvin})$, molar heat capacity $C_m = C/n$ $(\si{\joule\per\mole\per\kelvin})$, specific internal energy $u = U/m$, specific enthalpy $h = H/m$, specific entropy $s = S/m$, chemical potential $\mu_i = (\partial G/\partial n_i)_{T,P,n_{j\neq i}}$ $(\si{\joule\per\mole})$, thermal conductivity $k$ $(\si{\watt\per\meter\per\kelvin})$, dynamic viscosity $\mu$ $(\si{\pascal\second})$, emissivity $\varepsilon$, mole fraction $x_i = n_i/n_{\mathrm{tot}}$, and vapor quality $x = m_g/m_{\mathrm{tot}}$.
  \item \textbf{Universal Ratio Identity:} The ratio of \emph{any two extensive properties} of a homogeneous single phase is \textbf{always intensive} (e.g., $\rho = m/V$, $M = m/n$, $c = C/m$, $P = \tfrac{2}{3}U_{\mathrm{trans}}/V$). Conversely, the product of an extensive property and an intensive property is extensive ($U = mu$, $H = nH_m$).
\end{itemize}

\subsection{State Functions vs.\ Path Functions, Equilibrium vs.\ Steady State \& Quasistatic vs.\ Reversible}
\begin{itemize}
  \item \textbf{State Functions vs.\ Path Functions:} A \emph{state function} ($P, T, V, U, H, S, G, F, \rho$) is an equilibrium property whose change $\Delta X = X_2 - X_1$ depends \textbf{only on the endpoints $(1, 2)$} and satisfies $\Delta X_{\mathrm{cycle}} = 0$ around any closed loop. A \emph{path function} (heat $Q$, work $W_{\mathrm{by}}$, entropy generation $S_{\mathrm{gen}}$) describes an energy or entropy transfer \emph{during a process} and depends on the specific path taken between states $1$ and $2$; around a cycle, $Q_{\mathrm{net}} = W_{\mathrm{net,by}} \neq 0$ \src{S13,S16}.
  \item \textbf{Four Pillars of Thermodynamic Equilibrium vs.\ Steady State:} True thermodynamic equilibrium requires simultaneous (1)~\emph{thermal equilibrium} ($\nabla T = 0$, no heat flux), (2)~\emph{mechanical equilibrium} ($\nabla P = 0$ except gravity), (3)~\emph{phase equilibrium} ($\mu_i^{(\alpha)} = \mu_i^{(\beta)}$ across phases), and (4)~\emph{chemical equilibrium} ($\Delta_r G = \sum \nu_i \mu_i = 0$). A metal bar held between a $\SI{100}{\celsius}$ boiling bath and a $\SI{0}{\celsius}$ ice bath has a time-independent temperature profile ($\partial T/\partial t = 0$), so it is in \textbf{steady state}, NOT thermal equilibrium (it conducts a steady heat rate $\dot{Q} > 0$ and continuously generates entropy $\dot{S}_{\mathrm{gen}} = \dot{Q}(1/T_C - 1/T_H) > 0$) \src{S12,S16}.
  \item \textbf{Quasistatic vs.\ Reversible Trap:} A \emph{quasistatic} process proceeds slowly enough ($t_{\mathrm{proc}} \gg t_{\mathrm{relax}}$) that the system passes through a locus of well-defined internal equilibrium states. A \emph{reversible} process is quasistatic \textbf{and} free of all dissipative entropy generation ($S_{\mathrm{gen}} = 0$: no friction, no unrestrained expansion, no finite $\Delta T$ heat transfer, no mixing). Every reversible process is quasistatic, but a quasistatic process with friction (e.g., slowly pushing a rough piston) or across a finite temperature difference is \textbf{irreversible} ($S_{\mathrm{gen}} > 0$) \src{S13,S16}.
  \item \textbf{Zeroth Law of Thermodynamics (formulated by R.~H.~Fowler, 1931):} \emph{If body $A$ is in thermal equilibrium with body $C$, and body $B$ is separately in thermal equilibrium with body $C$, then $A$ and $B$ are in thermal equilibrium with each other} ($T_A = T_C \land T_B = T_C \implies T_A = T_B$) \src{S13}. This transitive equivalence relation defines \textbf{temperature} as the unique intensive property shared by all systems in mutual thermal equilibrium and validates using body $C$ as a \textbf{thermometer}.
\end{itemize}

\section{Thermometry \& Complete Eight-Scale Temperature Conversion System}
\label{sec:03.2}
\idxtopic{thermometers!constant-volume gas thermometer}
\idxtopic{temperature scales!Kelvin, Celsius, Fahrenheit, Rankine}
\idxtopic{temperature scales!Reaumur, Delisle, Newton, Romer}
\idxtopic{temperature scales!intersection points (-40, 574.59, +60)}

\subsection{Operational Thermometry \& Constant-Volume Ideal-Gas Thermometer}
Any physical property $X$ that varies monotonically with hotness (mercury/alcohol column length $L$, gas pressure $P$ at fixed volume, platinum RTD resistance $R(T) = R_0(1 + \alpha_R T_C)$, thermocouple Seebeck voltage) defines an empirical temperature scale. Because real materials expand nonlinearly at high/low temperatures, different liquid-in-glass thermometers calibrated at $\SI{0}{\celsius}$ and $\SI{100}{\celsius}$ disagree slightly at $\SI{50}{\celsius}$. However, all low-density gases inside a \textbf{constant-volume gas thermometer} converge to the exact same absolute thermodynamic temperature as reference pressure $P_{\mathrm{tp}} \to 0$ \src{S13}:
\begin{equation}\label{eq:03:gastherm}
  T(\si{\kelvin}) = (\SI{273.16}{\kelvin}) \lim_{P_{\mathrm{tp}}\to 0} \frac{P}{P_{\mathrm{tp}}},
\end{equation}
where $P_{\mathrm{tp}}$ is the gas pressure when the bulb is immersed in a triple-point cell of pure water ($T_{\mathrm{tp}} = \SI{273.16}{\kelvin} = \SI{0.01}{\celsius}$, $P_{\mathrm{tp,H_2O}} = \SI{611.657}{\pascal}$).

\subsection{Master Conversion Table \& Alignment Ladder for All 8 Competition Scales}
Science Olympiad written exams regularly test not only the four modern scales ($\si{\kelvin}, \si{\celsius}, ^\circ\mathrm{F}, ^\circ\mathrm{R}$) but also the four classical historical scales (\textbf{R\'eaumur, Delisle, Newton, R\o mer}) \src{S25,S26}. Every linear scale $S$ with water freezing point $S_{\mathrm{ice}}$ ($\SI{0}{\celsius}$) and normal boiling point $S_{\mathrm{boil}}$ ($\SI{100}{\celsius}$) obeys the universal algebraic ratio (\fref{fig:03:scales}, \tref{tab:03:scales}):
\begin{equation}\label{eq:03:universalscale}
  \frac{T_S - S_{\mathrm{ice}}}{S_{\mathrm{boil}} - S_{\mathrm{ice}}} = \frac{T_C}{100} = \frac{T_K - 273.15}{100} = \frac{T_F - 32}{180} = \frac{T_R - 491.67}{180} = \frac{T_{\mathrm{R\acute{e}}}}{80} = \frac{150 - T_{\mathrm{De}}}{150} = \frac{T_N}{33} = \frac{T_{\text{R\o}} - 7.5}{52.5}.
\end{equation}

\begin{center}
\captionof{table}{Complete conversion matrix, fixed points, and interval ratios ($\Delta T_S / \Delta T_C$) for all 8 competition temperature scales \src{S07,S08,S13}.}\label{tab:03:scales}
\footnotesize\setlength{\tabcolsep}{3.2pt}\renewcommand{\arraystretch}{1.15}
\begin{tabularx}{\linewidth}{@{}l c r r r r r l l c@{}}\toprule
\textbf{Scale (Year)} & \textbf{Sym.} & \textbf{Abs.\ Zero} & \textbf{Ice Pt.} & \textbf{Triple Pt.} & \textbf{Body} & \textbf{Steam Pt.} & \textbf{From Celsius ($\boldsymbol{T_C}$)} & \textbf{To Celsius ($\boldsymbol{T_C}$)} & $\boldsymbol{\Delta T_S/\Delta T_C}$ \\\midrule
\textbf{Kelvin} (1848) & $\si{\kelvin}$ & $0.00$ & $273.15$ & $273.16$ & $310.15$ & $373.15$ & $T_K = T_C + 273.15$ & $T_C = T_K - 273.15$ & $+1.000$ \\
\textbf{Celsius} (1742) & $\si{\celsius}$ & $-273.15$ & $0.00$ & $0.01$ & $37.00$ & $100.00$ & $T_C = T_C$ & $T_C = T_C$ & $+1.000$ \\
\textbf{Fahrenheit} (1724) & $^\circ\mathrm{F}$ & $-459.67$ & $32.00$ & $32.018$ & $98.60$ & $212.00$ & $T_F = \tfrac{9}{5}T_C + 32$ & $T_C = \tfrac{5}{9}(T_F - 32)$ & $+1.800$ \\
\textbf{Rankine} (1859) & $^\circ\mathrm{R}$ & $0.00$ & $491.67$ & $491.688$ & $558.27$ & $671.67$ & $T_R = \tfrac{9}{5}T_K = T_F + 459.67$ & $T_C = \tfrac{5}{9}T_R - 273.15$ & $+1.800$ \\
\textbf{R\'eaumur} (1730) & $^\circ\mathrm{R\acute{e}}$ & $-218.52$ & $0.00$ & $0.008$ & $29.60$ & $80.00$ & $T_{\mathrm{R\acute{e}}} = \tfrac{4}{5}T_C = 0.8\,T_C$ & $T_C = \tfrac{5}{4}T_{\mathrm{R\acute{e}}} = 1.25\,T_{\mathrm{R\acute{e}}}$ & $+0.800$ \\
\textbf{Delisle} (1732, \emph{inv.}) & $^\circ\mathrm{De}$ & $+559.725$ & $150.00$ & $149.985$ & $94.50$ & $0.00$ & $T_{\mathrm{De}} = \tfrac{3}{2}(100 - T_C)$ & $T_C = 100 - \tfrac{2}{3}T_{\mathrm{De}}$ & $\mathbf{-1.500}$ \\
\textbf{Newton} (1701) & $^\circ\mathrm{N}$ & $-90.1395$ & $0.00$ & $0.0033$ & $12.21$ & $33.00$ & $T_N = \tfrac{33}{100}T_C = 0.33\,T_C$ & $T_C = \tfrac{100}{33}T_N$ & $+0.330$ \\
\textbf{R\o mer} (1701) & $^\circ\text{R\o}$ & $-135.904$ & $7.50$ & $7.50525$ & $26.925$ & $60.00$ & $T_{\text{R\o}} = \tfrac{21}{40}T_C + 7.5$ & $T_C = \tfrac{40}{21}(T_{\text{R\o}} - 7.5)$ & $+0.525$ \\\bottomrule
\end{tabularx}
\end{center}

\begin{center}
\begin{tikzpicture}[>=Stealth,font=\scriptsize,xscale=1.02,yscale=0.68]
  % Horizontal alignment bars for 6 key fixed/crossover temperatures
  \foreach \y/\lbl/\tk/\tc/\tf/\tr/\tre/\tde/\tn/\tro in {
    0.0/{Absolute Zero}/0.00/-273.15/-459.67/0.00/-218.52/+559.73/-90.14/-135.90,
    1.1/{$^\circ\mathrm{C}=^\circ\mathrm{F}$ Crossover}/233.15/-40.00/-40.00/419.67/-32.00/+210.00/-13.20/-13.50,
    2.2/{$\mathrm{H_2O}$ Ice / Triple Pt.}/273.15/0.00/32.00/491.67/0.00/+150.00/0.00/+7.50,
    3.3/{Human Body ($37\,^\circ\mathrm{C}$)}/310.15/37.00/98.60/558.27/29.60/+94.50/12.21/+26.93,
    4.4/{$^\circ\mathrm{C}=^\circ\mathrm{De}$ Crossover}/333.15/60.00/140.00/599.67/48.00/+60.00/19.80/+39.00,
    5.5/{$\mathrm{H_2O}$ Steam Pt.\ ($1\text{ atm}$)}/373.15/100.00/212.00/671.67/80.00/0.00/33.00/+60.00
  }{
    \draw[gray!50,dashed] (2.8,\y) -- (16.8,\y);
    \node[anchor=east,font=\sffamily\bfseries\scriptsize,navy] at (2.7,\y) {\lbl};
    \node[fill=white,inner sep=1pt] at (3.6,\y) {\tk};
    \node[fill=white,inner sep=1pt] at (5.5,\y) {\tc};
    \node[fill=white,inner sep=1pt] at (7.4,\y) {\tf};
    \node[fill=white,inner sep=1pt] at (9.3,\y) {\tr};
    \node[fill=white,inner sep=1pt] at (11.2,\y) {\tre};
    \node[fill=white,inner sep=1pt] at (13.1,\y) {\tde};
    \node[fill=white,inner sep=1pt] at (14.9,\y) {\tn};
    \node[fill=white,inner sep=1pt] at (16.6,\y) {\tro};
  }
  % Vertical thermometer stems
  \foreach \x/\hdr in {3.6/{$\mathrm{K}$},5.5/{$^\circ\mathrm{C}$},7.4/{$^\circ\mathrm{F}$},9.3/{$^\circ\mathrm{R}$},11.2/{$^\circ\mathrm{R\acute{e}}$},13.1/{$^\circ\mathrm{De}$ (inv)},14.9/{$^\circ\mathrm{N}$},16.6/{$^\circ\text{R\o}$}}{
    \draw[thick,navy] (\x,-0.35) -- (\x,5.9);
    \node[above,font=\sffamily\bfseries\scriptsize,fill=boxbg,draw=navy,rounded corners=1pt,inner sep=2pt] at (\x,5.95) {\hdr};
  }
\end{tikzpicture}
\captionof{figure}{Eight-scale thermometric alignment ladder across absolute zero, key dual-scale intersection points ($-40^\circ$ and $+60^\circ$), and standard water fixed points. Note that Delisle ($^\circ\mathrm{De}$) increases \emph{downward} from boiling ($0\,^\circ\mathrm{De}$) to freezing ($150\,^\circ\mathrm{De}$) \src{S08,S13}.}\label{fig:03:scales}
\idxfig{temperature scales: eight-scale thermometric alignment ladder}
\end{center}

\begin{shortcut}[Six High-Yield Equal-Reading Crossover Temperatures]
\begin{itemize}
  \item \textbf{Celsius $=$ Fahrenheit:} $x = 1.8x + 32 \implies \mathbf{-40.0\,^\circ\mathrm{C} = -40.0\,^\circ\mathrm{F}}$ ($\SI{233.15}{\kelvin} = 419.67\,^\circ\mathrm{R}$).
  \item \textbf{Kelvin $=$ Fahrenheit:} $T_C + 273.15 = 1.8T_C + 32 \implies T_C = \SI{301.4375}{\celsius} \implies \mathbf{\SI{574.5875}{\kelvin} = 574.5875\,^\circ\mathrm{F}}$.
  \item \textbf{Celsius $=$ Delisle:} $x = 150 - 1.5x \implies 2.5x = 150 \implies \mathbf{+60.0\,^\circ\mathrm{C} = +60.0\,^\circ\mathrm{De}}$ ($\SI{333.15}{\kelvin} = 140.0\,^\circ\mathrm{F}$).
  \item \textbf{Fahrenheit $=$ Delisle:} $1.8T_C + 32 = 150 - 1.5T_C \implies T_C = \SI{35.7576}{\celsius} \implies \mathbf{+96.3636\,^\circ\mathrm{F} = +96.3636\,^\circ\mathrm{De}}$.
  \item \textbf{Celsius $=$ R\o mer:} $x = 0.525x + 7.5 \implies 0.475x = 7.5 \implies \mathbf{-15.7895\,^\circ\mathrm{C} = -15.7895\,^\circ\text{R\o}}$.
  \item \textbf{Celsius $=$ R\'eaumur $=$ Newton:} $\mathbf{0.0\,^\circ\mathrm{C} = 0.0\,^\circ\mathrm{R\acute{e}} = 0.0\,^\circ\mathrm{N}}$ (all three anchor zero at the ice point!).
\end{itemize}
\end{shortcut}

\section{Thermal Expansion of Solids, Liquids \& Gases}
\label{sec:03.3}
\idxtopic{thermal expansion!linear, superficial (2*alpha), volumetric (3*alpha)}
\idxtopic{thermal expansion!holes and cavities photographic enlargement}
\idxtopic{bimetallic strip!radius of curvature derivation}
\idxtopic{water 4 C density maximum anomaly!top-down lake freezing}

\subsection{Algebraic Binomial Derivations of 1D, 2D, and 3D Expansion}
At constant pressure, warming a solid or liquid increases the anharmonic average interatomic spacing \src{S13}:
\begin{enumerate}
  \item \textbf{Linear Expansion (1D rods, wires, perimeters, diameters):} $\Delta L = L - L_0 = \alpha_L L_0 \Delta T \implies L(T) = L_0(1 + \alpha_L\Delta T)$.
  \item \textbf{Superficial / Area Expansion (2D isotropic plates, circular holes, surface shells):} For an isotropic plate of dimensions $L_1 \times L_2$, each edge expands by $(1 + \alpha_L\Delta T)$:
    \[
      A(T) = A_0(1 + \alpha_L\Delta T)^2 = A_0\!\left[1 + 2\alpha_L\Delta T + (\alpha_L\Delta T)^2\right] \approx A_0(1 + 2\alpha_L\Delta T) \implies \boxed{\Delta A \approx 2\alpha_L A_0 \Delta T \quad (\alpha_A = 2\alpha_L)}.
    \]
  \item \textbf{Volumetric Expansion (3D blocks, hollow cavities, liquids):} For an isotropic solid $V(T) = V_0(1 + \alpha_L\Delta T)^3$:
    \[
      V(T) = V_0\!\left[1 + 3\alpha_L\Delta T + 3(\alpha_L\Delta T)^2 + (\alpha_L\Delta T)^3\right] \approx V_0(1 + 3\alpha_L\Delta T) \implies \boxed{\Delta V = \beta_V V_0 \Delta T, \quad \beta_{V,\mathrm{isotropic}} \approx 3\alpha_L}.
    \]
    For an \emph{anisotropic crystal} with principal linear expansion coefficients $(\alpha_x, \alpha_y, \alpha_z)$, $\beta_V = \alpha_x + \alpha_y + \alpha_z$. For an \emph{ideal gas} at constant pressure ($V = nRT/P$), $\beta_{V,\mathrm{ideal\ gas}} = \dfrac{1}{V}\dfrac{\Delta V}{\Delta T} = \dfrac{1}{T}$ (in kelvin: $\beta_V = 1/\SI{293.15}{\kelvin} = \SI{3.411e-3}{\per\kelvin}$ at $\SI{20}{\celsius}$) \src{S13}.
  \item \textbf{Holes, Annuli \& Hollow Cavities (Photographic Enlargement Theorem):} Every linear dimension of a uniform body---including the inner diameter of a washer or the interior volume of a hollow glass flask---expands \textbf{at the exact same rate as if the cavity were filled with the solid material itself}: $\Delta D_{\mathrm{hole}} = \alpha_{L,\mathrm{solid}} D_0 \Delta T$ and $\Delta V_{\mathrm{flask}} = 3\alpha_{L,\mathrm{glass}} V_0 \Delta T$ \src{S13}.
  \item \textbf{Apparent Liquid Overflow \& Density Variation:} When a vessel of volume $V_0$ and linear coefficient $\alpha_{\mathrm{ves}}$ is filled to the brim with a liquid of volumetric coefficient $\beta_{\mathrm{liq}}$ and heated by $\Delta T$:
    \begin{equation}\label{eq:03:overflow}
      \Delta V_{\mathrm{overflow}} = \Delta V_{\mathrm{liq}} - \Delta V_{\mathrm{ves}} = \left(\beta_{\mathrm{liq}} - 3\alpha_{\mathrm{ves}}\right) V_0 \Delta T, \qquad
      \rho(T) = \frac{m}{V(T)} = \frac{\rho_0}{1 + \beta_V\Delta T} \approx \rho_0(1 - \beta_V\Delta T).
    \end{equation}
\end{enumerate}

\subsection{Bimetallic Strips, Thermal Stress, Clock/Ruler Errors \& Water's $\SI{4}{\celsius}$ Anomaly}
\begin{itemize}
  \item \textbf{Bimetallic Strip Radius of Curvature ($R_{\mathrm{curv}}$):} Two metal strips of thickness $d$ (or center-to-center distance $d$) and expansion coefficients $\alpha_2 > \alpha_1$ are bonded at $T_0$ where both have length $L_0$ (\fref{fig:03:bimetallic}b). When heated by $\Delta T > 0$, strip~2 stretches more than strip~1, forcing the composite beam to bend into a circular arc of subtended angle $\theta$ (in radians) with strip~2 on the \textbf{outer (convex)} arc ($R_{\mathrm{curv}} + d/2$) and strip~1 on the \textbf{inner (concave)} arc ($R_{\mathrm{curv}} - d/2$):
    \[
      L_1 = \left(R_{\mathrm{curv}} - \tfrac{d}{2}\right)\theta = L_0(1 + \alpha_1\Delta T), \qquad
      L_2 = \left(R_{\mathrm{curv}} + \tfrac{d}{2}\right)\theta = L_0(1 + \alpha_2\Delta T).
    \]
    Subtracting gives $d\,\theta = L_0(\alpha_2 - \alpha_1)\Delta T$, while averaging gives $R_{\mathrm{curv}}\theta = L_0\left[1 + \tfrac{1}{2}(\alpha_1+\alpha_2)\Delta T\right] \approx L_0$. Dividing eliminates $\theta$ and $L_0$ by pure algebra \src{S26}:
    \begin{equation}\label{eq:03:bimetallic}
      \boxed{R_{\mathrm{curv}} = \frac{d\left[1 + \frac{1}{2}(\alpha_1 + \alpha_2)\Delta T\right]}{(\alpha_2 - \alpha_1)\Delta T} \approx \frac{d}{(\alpha_2 - \alpha_1)\Delta T}} \quad (\text{bends toward lower-}\alpha\text{ metal when heated, }\Delta T > 0).
    \end{equation}
  \item \textbf{Thermal Stress in a Clamped Rod \& Buckling:} If a rod of Young's modulus $E$ ($\sigma = E\,\Delta L_{\mathrm{mech}}/L_0$) is held between rigid walls so $\Delta L_{\mathrm{tot}} = \Delta L_{\mathrm{th}} + \Delta L_{\mathrm{mech}} = 0$, heating by $\Delta T$ induces compressive stress $\sigma = F/A = -E\alpha_L\Delta T$ \src{S13}.
  \item \textbf{Metal Scale / Tape-Measure \& Pendulum Clock Corrections:} A metal ruler calibrated at $T_0$ has enlarged tick spacing at $T > T_0$, so it \emph{under-reads} true length $L_{\mathrm{true}} = L_{\mathrm{read}}\left[1 + \alpha_{\mathrm{ruler}}(T - T_0)\right]$. A pendulum clock ($P_{\mathrm{pend}} = 2\pi\sqrt{L/g} \propto L^{1/2}$) has fractional period shift $\Delta P_{\mathrm{pend}}/P_{\mathrm{pend}} \approx +\tfrac{1}{2}\alpha_L\Delta T$: it runs \textbf{slow} when warm, losing $\Delta t \approx \tfrac{1}{2}\alpha_L\Delta T \times \SI{86400}{\second/day}$.
  \item \textbf{Water's $\SI{3.984}{\celsius}$ ($\approx\SI{4.0}{\celsius}$) Density Maximum Anomaly:} Between $\SI{0.00}{\celsius}$ ($\rho = \SI{999.84}{\kilogram\per\cubic\meter}$) and $\SI{3.984}{\celsius}$ ($\rho_{\max} = \SI{999.97}{\kilogram\per\cubic\meter}$), liquid water has a \textbf{negative} volumetric expansion coefficient ($\beta_V = \SI{-68e-6}{\per\kelvin}$ at $\SI{0.01}{\celsius}$, crossing zero at $\SI{3.984}{\celsius}$) because open tetrahedral hydrogen-bonded ice-like clusters collapse slightly upon warming (\fref{fig:03:bimetallic}a) \src{S09,S13}. In winter, cooling lake surface water sinks until the entire water column reaches $\SI{4}{\celsius}$; further cooling below $\SI{4}{\celsius}$ makes the surface water \emph{lighter} ($\beta_V < 0$), so $\SIrange{0}{4}{\celsius}$ water and $\SI{0}{\celsius}$ ice ($\rho_{\mathrm{ice}} = \SI{916.7}{\kilogram\per\cubic\meter}$) stay stratified at the top, insulating aquatic life at $\SI{4}{\celsius}$ below!
\end{itemize}

\begin{center}
\begin{tikzpicture}[>=Stealth,font=\scriptsize]
  % Left panel: Water 4C density curve from NIST IAPWS-95
  \begin{scope}
    \begin{axis}[
      width=7.8cm, height=4.1cm,
      xlabel={Temperature $T$ ($\si{\celsius}$)},
      ylabel={Liquid Water Density $\rho$ ($\si{\kilogram\per\cubic\meter}$)},
      xmin=0, xmax=20, ymin=998.0, ymax=1000.1,
      xtick={0,4,8,12,16,20},
      ytick={998.2,998.6,999.0,999.4,999.8,1000.0},
      grid=major, grid style={gray!30},
      title={\textbf{(a) NIST IAPWS-95 Water Density Maximum at $\SI{3.984}{\celsius}$}},
      label style={font=\scriptsize}, tick label style={font=\scriptsize}, title style={font=\scriptsize}
    ]
      \addplot[very thick,navy,smooth,mark=*,mark size=1.2pt] coordinates {
        (0.01,999.84) (2.0,999.94) (3.984,999.97) (5.0,999.96) (7.0,999.90)
        (10.0,999.70) (12.0,999.50) (15.0,999.10) (17.0,998.78) (20.0,998.21)
      };
      \draw[dashed,rust,thick] (axis cs:3.984,998.0) -- (axis cs:3.984,999.97);
      \node[anchor=north west,font=\scriptsize,rust] at (axis cs:4.2,999.45) {$\rho_{\max}=\SI{999.97}{\kilogram\per\cubic\meter}$ ($\beta_V=0$)};
      \node[anchor=south west,font=\scriptsize,teal] at (axis cs:0.2,998.15) {$\beta_V < 0$ ($0\text{--}4\,^\circ\mathrm{C}$)};
    \end{axis}
  \end{scope}

  % Right panel: Bimetallic strip curvature
  \begin{scope}[xshift=8.6cm,yshift=0.25cm]
    \node[font=\sffamily\bfseries\scriptsize] at (4.1,3.35) {(b) Bimetallic Strip Bending Geometry ($\alpha_2 > \alpha_1$)};
    % Cold straight state
    \draw[fill=boxbg,draw=navy,thick] (0.3,2.55) rectangle (3.5,2.8);
    \draw[fill=gray!45,draw=navy,thick] (0.3,2.3) rectangle (3.5,2.55);
    \node[right,font=\scriptsize] at (3.55,2.68) {Metal 2 ($\alpha_2$ higher, e.g., Brass $\num{19e-6}$)};
    \node[right,font=\scriptsize] at (3.55,2.42) {Metal 1 ($\alpha_1$ lower, e.g., Invar $\num{1.2e-6}$)};
    \node[above,font=\bfseries\scriptsize,navy] at (1.9,2.82) {At $T_0$ ($\Delta T = 0$): flat, $L_1 = L_2 = L_0$};
    % Heated curved arc
    \draw[thick,navy,fill=boxbg] (0.5,0.85) arc (140:40:2.3cm) -- ++(-0.18,-0.15) arc (40:140:2.06cm) -- cycle;
    \draw[thick,navy,fill=gray!45] (0.68,0.70) arc (140:40:2.06cm) -- ++(-0.18,-0.15) arc (40:140:1.82cm) -- cycle;
    \draw[dashed,rust,->] (2.26,-0.55) -- (2.26,1.45) node[midway,right,font=\bfseries] {$R_{\mathrm{curv}} \approx \dfrac{d}{(\alpha_2-\alpha_1)\Delta T}$};
    \node[above,font=\bfseries\scriptsize,rust] at (2.3,1.72) {Heated ($\Delta T > 0$): outer arc = Metal 2 ($\alpha_2$)};
  \end{scope}
\end{tikzpicture}
\captionof{figure}{(a)~Anomalous density maximum of liquid water at $T = \SI{3.984}{\celsius}$ ($\rho_{\max} = \SI{999.97}{\kilogram\per\cubic\meter}$ at $\SI{1}{atm}$) where $\beta_V$ changes sign from negative to positive \src{S09}. (b)~Bimetallic strip bending toward the lower-expansion metal ($\alpha_1$) upon heating ($\Delta T > 0$) with radius $R_{\mathrm{curv}} \approx d/[(\alpha_2-\alpha_1)\Delta T]$ \src{S13,S26}.}\label{fig:03:bimetallic}
\idxfig{water 4 C density anomaly and bimetallic strip curvature}
\end{center}

\section{Historical Perspective of the Exact 8 Named Scientists (Rule \texttt{5.c.v})}
\label{sec:03.4}
\idxtopic{Rule 5.c.v!exact 8 named scientists}
\idxtopic{Kelvin (William Thomson)!1848 scale, 1851 dynamical theory, Joule-Thomson}
\idxtopic{Black, Joseph!latent heat (1761), specific heat, fixed air CO2}
\idxtopic{Joule, James Prescott!paddle-wheel mechanical equivalent, I^2 R, free expansion}
\idxtopic{Carnot, Sadi!1824 Motive Power of Fire, Carnot cycle and theorem}
\idxtopic{Planck, Max!1900 blackbody quantum E=h*nu, Third Law S(0)=0}
\idxtopic{Clausius, Rudolf!1850 1st/2nd Law, 1865 coined entropy, mean free path}
\idxtopic{Boltzmann, Ludwig!1877 S = k_B ln W, transport equation, Stefan-Boltzmann proof}
\idxtopic{Maxwell, James Clerk!1860 speed distribution, 4 Maxwell relations, Maxwell's demon}

\begin{rulebox}[2027 Rule 5.c.v Mandatory Historical Syllabus \src{S01}]
Rule \texttt{5.I.c} requires \textbf{at least 3 written questions} at every tournament level (Regionals, States, Nationals) from Clause \texttt{5.c.v}: \emph{``Historical perspective of contributions of \textbf{Kelvin, Joseph Black, Joule, Carnot, Planck, Clausius, Boltzmann, and Maxwell}.''} The timeline (\fref{fig:03:history}), deep profiles (\tref{tab:03:scientists}), and decoy disambiguation matrix (\tref{tab:03:historytraps}) cover every competition question pattern on these 8 scientists \src{S23,S25}.
\end{rulebox}

\begin{center}
\begin{tikzpicture}[>=Stealth,font=\scriptsize,x=0.071cm,y=0.68cm]
  % Time axis from 1725 (x=0) to 1955 (x=230)
  \draw[->,ultra thick,navy] (0,0) -- (230,0) node[right,font=\bfseries] {Year};
  \foreach \yr/\x in {1725/0,1750/25,1775/50,1800/75,1825/100,1850/125,1875/150,1900/175,1925/200,1950/225}{
    \draw[thick,navy] (\x,-0.15) -- (\x,0.15) node[below=4pt,font=\bfseries\scriptsize] {\yr};
  }
  % Lifespan bars of the 8 Named Scientists (x = Year - 1725)
  \draw[very thick,teal,fill=boxbg] (3,4.2) rectangle (74,4.65) node[midway,font=\bfseries] {2. Joseph Black (1728--1799)};
  \node[right,font=\scriptsize,rust] at (75,4.42) {1754 $\mathrm{CO_2}$; 1761 Latent \& Specific Heat};

  \draw[very thick,rust,fill=boxbg] (71,3.3) rectangle (107,3.75) node[pos=0.5,above=4pt,font=\bfseries] {4. Sadi Carnot (1796--1832): 1824 \emph{Motive Power of Fire} ($\eta_C$)};

  \draw[very thick,navy,fill=boxbg] (93,2.4) rectangle (164,2.85) node[midway,font=\bfseries] {3. James Prescott Joule (1818--1889)};
  \node[right,font=\scriptsize,navy] at (165,2.62) {1840 $I^2R$; 1845--50 Paddle Wheel $J$};

  \draw[very thick,teal,fill=boxbg] (97,1.5) rectangle (163,1.95) node[midway,font=\bfseries] {6. Rudolf Clausius (1822--1888)};
  \node[right,font=\scriptsize,teal] at (164,1.72) {1850 1st/2nd Law; 1865 Coined \emph{Entropy} $S$};

  \draw[very thick,amber,fill=boxbg] (99,0.6) rectangle (182,1.05) node[midway,font=\bfseries] {1. Lord Kelvin (1824--1907)};
  \node[right,font=\scriptsize,amber] at (183,0.82) {1848 Abs.\ Scale; 1852 $\mu_{\mathrm{JT}}$};

  \draw[very thick,rust,fill=boxbg] (106,-1.15) rectangle (154,-0.7) node[midway,font=\bfseries] {8. J.\ C.\ Maxwell (1831--79)};
  \node[right,font=\scriptsize,rust] at (155,-0.92) {1860 $f(v)$; 1871 4 Relations \& Demon};

  \draw[very thick,navy,fill=boxbg] (119,-2.05) rectangle (181,-1.6) node[midway,font=\bfseries] {7. L.\ Boltzmann (1844--1906)};
  \node[right,font=\scriptsize,navy] at (182,-1.82) {1877 $S = k_B\ln\Omega$; 1884 $\sigma T^4$};

  \draw[very thick,teal,fill=boxbg] (133,-2.95) rectangle (222,-2.5) node[midway,font=\bfseries] {5. Max Planck (1858--1947)};
  \node[left,font=\scriptsize,teal] at (132,-2.72) {1900 Blackbody $E=h\nu$; 1911 3rd Law $S(0)=0$};
\end{tikzpicture}
\captionof{figure}{Chronological lifespan and discovery timeline of the exact 8 scientists specified in 2027 Rule \texttt{5.c.v} (1728--1947) \src{S01,S23}.}\label{fig:03:history}
\idxfig{timeline of the 8 named scientists in Rule 5.c.v}
\end{center}

\begin{center}
\captionof{table}{Authoritative historical reference profiles for the exact 8 scientists in 2027 Rule \texttt{5.c.v} \src{S01,S23}.}\label{tab:03:scientists}
\footnotesize\setlength{\tabcolsep}{3pt}\renewcommand{\arraystretch}{1.15}
\begin{tabularx}{\linewidth}{@{}>{\raggedright\arraybackslash}p{0.17\linewidth} >{\raggedright\arraybackslash}p{0.39\linewidth} >{\raggedright\arraybackslash}p{0.22\linewidth} >{\raggedright\arraybackslash}X@{}}\toprule
\textbf{Scientist (\texttt{5.c.v})} & \textbf{Key Experiments, Publications \& Years} & \textbf{Laws, Equations \& Units Named} & \textbf{Distinguishing Exam Traps} \\\midrule
\textbf{1. William Thomson (Lord Kelvin)} (1824--1907, Belfast / Glasgow, British) & \textbf{1848:} \emph{On an Absolute Thermometric Scale} based on Carnot cycles. \textbf{1851:} \emph{On the Dynamical Theory of Heat} reconciling Carnot with Joule. \textbf{1852--62:} Porous-plug throttling with Joule. \textbf{1854:} Thermoelectric Thomson effect \& $\Pi = \alpha T$. & SI unit \textbf{kelvin} ($\si{\kelvin}$); Kelvin--Planck 2nd Law statement; Joule--Thomson $\mu_{\mathrm{JT}} = (\partial T/\partial P)_H$; Kelvin scale $T_C/T_H = Q_C/Q_H$. & Initially defended caloric theory at 1847 BAAS meeting before reconciling Carnot \& Joule in 1851. Coined \emph{``Maxwell's demon''} (1874). Title from River Kelvin at Univ.\ of Glasgow. \\\midrule
\textbf{2. Joseph Black} (1728--1799, Bordeaux / Glasgow / Edinburgh, Scottish) & \textbf{1754:} Isolated $\mathrm{CO_2}$ (\emph{``fixed air''}) gravimetrically from $\mathrm{MgCO_3/CaCO_3}$. \textbf{1757--60:} Discovered \textbf{specific heat} (mercury vs.\ water mixing). \textbf{1761--62:} Discovered \textbf{latent heat} of fusion (ice melting) \& vaporization; mentored/funded James Watt's separate condenser (1765). & Latent heat $Q = mL$ (coined \emph{latent heat}, 1761); specific heat \& method of mixtures $\sum m_i c_i \Delta T_i = 0$; heat vs.\ temperature distinction. & Taught discoveries in lectures (published posthumously 1803 by Robison). Do \textbf{not} confuse with Lavoisier \& Laplace (who built the 1780 \emph{ice calorimeter} using Black's $L_f$). \\\midrule
\textbf{3. James Prescott Joule} (1818--1889, Salford / Manchester, English) & \textbf{1840:} Ohmic heating $P = I^2R$. \textbf{1845:} Free expansion of air into vacuum in twin cylinders ($\Delta T = 0 \implies U_{\mathrm{ideal}} = U(T)$). \textbf{1845--50:} Falling-weight brass \textbf{paddle-wheel calorimeter} in water, mercury \& sperm oil ($J = \SI{772.69}{ft\cdot lbf/BTU}$, $<0.7\%$ from $\SI{4.184}{\joule/cal}$). & SI energy unit \textbf{joule} ($\si{\joule}$); Mechanical equivalent of heat $W = JQ$; Joule's 1st Law $Q = I^2Rt$; Joule's 2nd Law $(\partial U/\partial V)_T = 0$; Joule--Thomson effect. & Distinguish Joule free expansion (closed rigid chambers into vacuum, $\Delta U = 0$) from Joule--Thomson throttling (steady porous-plug flow, $\Delta H = 0$). Mayer proposed conservation in 1842; Joule proved it quantitatively. \\\midrule
\textbf{4. Nicolas L\'eonard Sadi Carnot} (1796--1832, Paris, French military engineer) & \textbf{1824:} Published his sole book \emph{R\'eflexions sur la puissance motrice du feu} (\emph{Reflections on the Motive Power of Fire}) at age 28, introducing the 4-step reversible Carnot cycle (2 isotherms + 2 adiabats) and proving $\eta_{\max}$ depends only on $T_H, T_C$, independent of working fluid. & \textbf{Carnot cycle}; \textbf{Carnot efficiency} $\eta_C = 1 - T_C/T_H$; \textbf{Carnot's theorem}; Carnot $\mathrm{COP}_R = T_C/(T_H-T_C)$ \& $\mathrm{COP}_{HP} = T_H/(T_H-T_C)$. & Used caloric waterfall analogy in 1824 ($Q_H = Q_C$), but his unpublished pre-1832 notes rejected caloric and estimated $J \approx \SI{3.7}{\joule/cal}$. Clapeyron (1834) first drew Carnot's cycle on a $P$--$V$ diagram. \\\midrule
\textbf{5. Max Planck} (1858--1947, Kiel / Munich / Berlin, German; 1918 Nobel) & \textbf{1879/1897:} PhD \& \emph{Treatise on Thermodynamics} (2nd Law engine statement). \textbf{Oct/Dec 14, 1900:} Derived \textbf{Planck blackbody spectrum} by quantizing oscillator energy $E = h\nu$, resolving the Rayleigh--Jeans UV catastrophe, and evaluated $h$ and $k_B$. \textbf{1911:} Absolute Third Law ($S \to 0$ as $T \to 0$ for perfect crystal). & \textbf{Planck constant} $h = \SI{6.62607015e-34}{\joule\second}$; \textbf{Planck's radiation law} $M_\lambda(\lambda,T)$; Kelvin--Planck 2nd Law; Planck 3rd Law ($\lim_{T\to 0}S_{\mathrm{cryst}} = 0$). & Planck first wrote Boltzmann's equation explicitly as $S = k\log W$ in Dec 1900 and named $k$ Boltzmann's constant! Walther Nernst discovered the Nernst Heat Theorem ($\Delta S \to 0$) in 1906; Planck set $S(0)=0$ in 1911. \\\midrule
\textbf{6. Rudolf Clausius} (1822--1888, K\"oslin / Zurich / Bonn, German) & \textbf{1850:} \emph{On the Moving Force of Heat}---first joint mathematical formulation of 1st Law ($dQ = dU + P\,dV$) and 2nd Law (refrigerator statement). \textbf{1857--58:} Kinetic theory rotational/vibrational modes \& \textbf{mean free path} $\lambda$. \textbf{1865:} Coined \textbf{Entropy} ($S$, $dS = \delta Q_{\mathrm{rev}}/T$) and $\oint \delta Q/T \leq 0$. & \textbf{Clausius statement} of 2nd Law; \textbf{Clausius inequality} $\oint \delta Q/T \leq 0$; macroscopic \textbf{entropy} $\Delta S = Q_{\mathrm{rev}}/T$; \textbf{Clausius--Clapeyron equation}; Virial theorem (1870). & Famous 1865 couplet: \emph{``1.~The energy of the universe is constant. 2.~The entropy of the universe tends to a maximum.''} Clausius introduced $\lambda$ with factor $\tfrac{3}{4}$; Maxwell corrected it to $1/\sqrt{2}$ in 1860. \\\midrule
\textbf{7. Ludwig Boltzmann} (1844--1906, Vienna / Graz, Austrian) & \textbf{1868--71:} Generalized Maxwell speeds with \textbf{Boltzmann factor} $e^{-E/(k_B T)}$ and proved general equipartition ($\tfrac{1}{2}k_B T$). \textbf{1872:} Boltzmann transport equation \& $H$-theorem ($dH/dt \leq 0$). \textbf{1877:} Statistical mechanics entropy $S = k_B\ln\Omega$. \textbf{1884:} Thermodynamic Carnot proof of Stefan's $M_e = \sigma T^4$ using radiation pressure $P = u/3$. & \textbf{Boltzmann constant} $k_B = \SI{1.380649e-23}{\joule\per\kelvin}$; \textbf{Boltzmann entropy} $S = k\log W$ (on Vienna tombstone); \textbf{Stefan--Boltzmann law} $\sigma T^4$; Maxwell--Boltzmann distribution. & Josef Stefan (Boltzmann's PhD advisor) found $M_e \propto T^4$ \emph{empirically} in 1879; Boltzmann proved it \emph{theoretically} in 1884. J.~W.~Gibbs coined ``statistical mechanics'' (1902) and defined $G = H - TS$. \\\midrule
\textbf{8. James Clerk Maxwell} (1831--1879, Edinburgh / London / Cambridge, Scottish) & \textbf{1860:} Derived \textbf{Maxwell speed distribution} $f(v)$ and predicted dilute gas viscosity $\mu$ is independent of pressure/density (verified experimentally with Katherine Maxwell in 1866). \textbf{1867/1871:} \emph{Theory of Heat}---derived the \textbf{Four Maxwell Relations} and invented the \textbf{Maxwell's Demon} thought experiment. & \textbf{Maxwell speed distribution} ($v_{\mathrm{mp}} < v_{\mathrm{avg}} < v_{\mathrm{rms}}$); \textbf{Four Maxwell Relations} from $dU, dH, dF, dG$; \textbf{Maxwell's Demon}; Maxwell equal-area rule (1875). & Maxwell's Demon proves the Second Law is \emph{statistical} (resolved via Landauer/Szilard information erasure entropy $k_B\ln 2$ per bit). Gas viscosity $\mu \sim \rho \bar{v}\lambda$ is independent of $P$ because $\rho \propto P$ and $\lambda \propto 1/P$ cancel! \\\bottomrule
\end{tabularx}
\end{center}

\begin{center}
\captionof{table}{High-frequency historical attribution traps: 2027 Rule \texttt{5.c.v} scientists vs.\ related thermodynamic pioneers \src{S23,S25}.}\label{tab:03:historytraps}
\footnotesize\setlength{\tabcolsep}{3.5pt}\renewcommand{\arraystretch}{1.13}
\begin{tabularx}{\linewidth}{@{}p{0.31\linewidth} p{0.25\linewidth} X@{}}\toprule
\textbf{Discovery / Milestone Question Cue} & \textbf{Correct Scientist} & \textbf{Common Exam Trap / How to Distinguish} \\\midrule
Coined \emph{latent heat} (1761), discovered specific heat \& isolated $\mathrm{CO_2}$ (\emph{fixed air}) & \textbf{Joseph Black} (\texttt{5.c.v}) & \textbf{Lavoisier \& Laplace} built the first \emph{ice calorimeter} (1780--83) and coined \emph{calorimeter} using Black's latent heat; Lavoisier named the \emph{caloric} fluid theory (1783). \\
1798 Munich cannon-boring demonstration that frictional heat is inexhaustible & \textbf{Count Rumford} (Benjamin Thompson) & Rumford attacked caloric \emph{qualitatively} via brass cannon boring; \textbf{Joule} (\texttt{5.c.v}) measured the exact \emph{quantitative} mechanical equivalent $J$ via paddle wheels (1845--50). \\
First philosophical conservation of energy \& $C_P - C_V = R$ gas estimate of $J$ (1842) & \textbf{Julius Robert von Mayer} & Mayer (1842) published the first conceptual statement and gas-expansion calculation; \textbf{Joule} (\texttt{5.c.v}) provided the precision experimental proof; \textbf{Helmholtz} (1847) formalized potential + kinetic conservation. \\
First drew Carnot's cycle on a $P$--$V$ indicator diagram \& derived $dP/dT = L/(T\Delta v)$ (1834) & \textbf{\'Emile Clapeyron} & \textbf{Carnot} (\texttt{5.c.v}, 1824) wrote verbally without $P$--$V$ graphs; Clapeyron (1834) graphed it; \textbf{Clausius} (\texttt{5.c.v}, 1850) corrected it for energy conservation and derived the ideal-vapor Clausius--Clapeyron form. \\
Empirical discovery of $M_e = \sigma T^4$ (1879) vs.\ thermodynamic Carnot derivation (1884) & \textbf{Josef Stefan} (1879 exp.) vs.\ \textbf{Ludwig Boltzmann} (\texttt{5.c.v}, 1884 theory) & Stefan deduced $T^4$ from Tyndall's platinum data (1879); his student \textbf{Boltzmann} (\texttt{5.c.v}) derived it theoretically from Maxwell's radiation pressure $P = u/3$ in 1884. \\
Wien displacement law $\lambda_{\max}T = b$ (1893) vs.\ quantum blackbody spectrum $E = h\nu$ (1900) & \textbf{Wilhelm Wien} (1893) vs.\ \textbf{Max Planck} (\texttt{5.c.v}, 1900) & Wien used classical Doppler/adiabatic cavity compression (1893); \textbf{Planck} (\texttt{5.c.v}) quantized oscillator energy $E=h\nu$ in 1900 (from which both Wien's $b$ and Stefan's $\sigma$ follow analytically). \\
Nernst Heat Theorem $\Delta S \to 0$ (1906) vs.\ perfect-crystal $S(0\,\si{\kelvin}) = 0$ (1911) & \textbf{Walther Nernst} (1906) vs.\ \textbf{Max Planck} (\texttt{5.c.v}, 1911) & Nernst formulated the Third Law in 1906 (1920 Chemistry Nobel); \textbf{Planck} (\texttt{5.c.v}) extended it in 1911 to fix the absolute entropy zero $S(0\,\si{\kelvin}) = 0$ for any pure perfect crystal. \\\bottomrule
\end{tabularx}
\end{center}

\section{Layered Worked Examples \& Diagnostics}
\label{sec:03.5}

\begin{workedexample}[Historical Temperature Scale Conversion: Delisle \& R\'eaumur to Kelvin]{Basic $\cdot$ Rule 5.c.i}{ex:03.1}
In 1735, a Russian meteorological station recorded a winter air temperature of $T_{\mathrm{De}} = 195.0\,^\circ\mathrm{De}$ on a Delisle thermometer, while a French chemist recorded an alcohol bath at $T_{\mathrm{R\acute{e}}} = -24.0\,^\circ\mathrm{R\acute{e}}$. Convert both readings to $\si{\celsius}$, $\si{\kelvin}$, and $^\circ\mathrm{F}$, and determine which is colder.

\textbf{Solution.}
For the Delisle reading ($T_{\mathrm{De}} = 195.0\,^\circ\mathrm{De}$, inverted scale where $150\,^\circ\mathrm{De} = \SI{0}{\celsius}$ and $0\,^\circ\mathrm{De} = \SI{100}{\celsius}$):
\[
  T_{C,1} = 100 - \frac{2}{3}T_{\mathrm{De}} = 100 - \frac{2}{3}(195.0) = 100 - 130.0 = \mathbf{-\SI{30.0}{\celsius}}, \qquad
  T_{K,1} = -30.0 + 273.15 = \mathbf{\SI{243.15}{\kelvin} \approx \SI{243.2}{\kelvin}},
\]
\[
  T_{F,1} = \frac{9}{5}(-30.0) + 32 = -54.0 + 32 = \mathbf{-22.0\,^\circ\mathrm{F}}.
\]
For the R\'eaumur reading ($T_{\mathrm{R\acute{e}}} = -24.0\,^\circ\mathrm{R\acute{e}}$, where $80\,^\circ\mathrm{R\acute{e}} = \SI{100}{\celsius}$):
\[
  T_{C,2} = \frac{5}{4}T_{\mathrm{R\acute{e}}} = 1.25(-24.0) = \mathbf{-\SI{30.0}{\celsius}}, \qquad
  T_{K,2} = \mathbf{\SI{243.15}{\kelvin}}, \qquad T_{F,2} = \mathbf{-22.0\,^\circ\mathrm{F}}.
\]
Both readings represent the \textbf{exact same temperature} ($-\SI{30.0}{\celsius} = \SI{243.15}{\kelvin} = -22.0\,^\circ\mathrm{F}$)!
\verifiedby{Universal ratio check (\ref{eq:03:universalscale}): $(150 - 195.0)/150 = -0.300$ and $-24.0/80 = -0.300 = T_C/100$ (\texttt{check\_sec02\_05.py}).}
\end{workedexample}

\begin{workedexample}[Shrink-Fit Interference Assembly of a Brass Ring onto a Steel Shaft]{Competition $\cdot$ Rule 5.c.i}{ex:03.2}
At $T_0 = \SI{20.0}{\celsius}$, a circular brass collar ($\alpha_{\mathrm{brass}} = \SI{19.0e-6}{\per\celsius}$) has an inner hole diameter of $D_{\mathrm{brass},0} = \SI{4.990}{\centi\meter}$, while a cylindrical steel shaft ($\alpha_{\mathrm{steel}} = \SI{12.0e-6}{\per\celsius}$) has an outer diameter of $D_{\mathrm{steel},0} = \SI{5.000}{\centi\meter}$. Find (a) the minimum temperature to which ONLY the brass collar must be heated so it slips over the $\SI{20.0}{\celsius}$ steel shaft, and (b) the common temperature to which BOTH the brass collar and steel shaft must be heated together so the collar slips onto the shaft.

\textbf{Solution.}
\textbf{(a)} Heating only the brass collar until $D_{\mathrm{brass}}(T) = D_{\mathrm{brass},0}(1 + \alpha_{\mathrm{brass}}\Delta T) = D_{\mathrm{steel},0} = \SI{5.000}{\centi\meter}$:
\[
  \Delta T = \frac{D_{\mathrm{steel},0} - D_{\mathrm{brass},0}}{\alpha_{\mathrm{brass}}\,D_{\mathrm{brass},0}} = \frac{\SI{0.010}{\centi\meter}}{(\SI{19.0e-6}{\per\celsius})(\SI{4.990}{\centi\meter})} = \SI{105.47}{\celsius} \implies T = 20.0 + 105.5 = \SI{125.5}{\celsius} \approx \mathbf{\SI{1.3e2}{\celsius}}.
\]
\textbf{(b)} Heating both together to $T = T_0 + \Delta T$ requires $D_{\mathrm{brass},0}(1 + \alpha_{\mathrm{brass}}\Delta T) = D_{\mathrm{steel},0}(1 + \alpha_{\mathrm{steel}}\Delta T)$:
\[
  \Delta T = \frac{D_{\mathrm{steel},0} - D_{\mathrm{brass},0}}{\alpha_{\mathrm{brass}} D_{\mathrm{brass},0} - \alpha_{\mathrm{steel}} D_{\mathrm{steel},0}} = \frac{0.010}{19.0\times 10^{-6}(4.990) - 12.0\times 10^{-6}(5.000)} = \frac{0.010}{\num{3.481e-5}} = \SI{287.27}{\celsius}
\]
so $T = \SI{307.3}{\celsius} \approx \mathbf{\SI{3.1e2}{\celsius}}$ (2 s.f.).
\textbf{Sig figs:} the gap $5.000 - 4.990 = \SI{0.010}{\centi\meter}$ has only 2 significant figures, so both temperatures are reported to 2 (full values kept in the working).
\verifiedby{At $T = \SI{307.27}{\celsius}$, $D_{\mathrm{brass}} = 4.990(1 + 19\times 10^{-6}\times 287.274) = \SI{5.01723}{\centi\meter} = D_{\mathrm{steel}}$ (\texttt{check\_sec02\_05.py}).}
\end{workedexample}

\begin{workedexample}[Apparent Overflow of Ethanol in a Pyrex Beaker]{Competition $\cdot$ Rule 5.c.i}{ex:03.3}
A Pyrex glass beaker ($\alpha_{L,\mathrm{Pyrex}} = \SI{3.0e-6}{\per\celsius}$) has an interior volume of exactly $V_0 = \SI{500.0}{\milli\liter}$ at $T_1 = \SI{15.0}{\celsius}$ and is filled to the brim with liquid ethanol ($\beta_{\mathrm{ethanol}} = \SI{1100e-6}{\per\celsius}$). When the assembly is warmed uniformly to $T_2 = \SI{45.0}{\celsius}$ ($\Delta T = \SI{30.0}{\celsius}$), how many milliliters of ethanol overflow?

\textbf{Solution.}
The volumetric expansion coefficient of isotropic Pyrex glass is $\beta_{\mathrm{Pyrex}} = 3\alpha_{L,\mathrm{Pyrex}} = 3(\SI{3.0e-6}{\per\celsius}) = \SI{9.0e-6}{\per\celsius}$. By (\ref{eq:03:overflow}):
\[
  \Delta V_{\mathrm{overflow}} = (\beta_{\mathrm{ethanol}} - 3\alpha_{L,\mathrm{Pyrex}}) V_0 \Delta T = (1100 - 9.0)\times 10^{-6}\,\si{\per\celsius} \times (\SI{500.0}{\milli\liter})(\SI{30.0}{\celsius}) = \mathbf{\SI{16.365}{\milli\liter} \approx \SI{16.4}{\milli\liter}}.
\]
\verifiedby{$\Delta V_{\mathrm{liq}} = \SI{16.500}{\milli\liter}$ minus $\Delta V_{\mathrm{glass}} = \SI{0.135}{\milli\liter} = \SI{16.365}{\milli\liter}$ (\texttt{check\_sec02\_05.py}).}
\end{workedexample}

\begin{workedexample}[Bimetallic Strip Curvature \& Tip Deflection (Brass--Invar)]{National $\cdot$ Rule 5.c.i}{ex:03.4}
A bimetallic thermostat cantilever of length $L_0 = \SI{8.00}{\centi\meter}$ is fabricated by bonding a $\SI{0.50}{\milli\meter}$-thick brass strip ($\alpha_2 = \SI{19.0e-6}{\per\celsius}$) to a $\SI{0.50}{\milli\meter}$-thick Invar strip ($\alpha_1 = \SI{1.2e-6}{\per\celsius}$) at $\SI{20.0}{\celsius}$. When heated to $\SI{120.0}{\celsius}$ ($\Delta T = \SI{100.0}{\celsius}$), compute (a) the center-to-center spacing $d$ and radius of curvature $R_{\mathrm{curv}}$, and (b) the transverse tip deflection $\delta_{\mathrm{tip}} \approx L_0^2 / (2R_{\mathrm{curv}})$ using the small-sagitta chord geometry $L_0^2 \approx 2R_{\mathrm{curv}}\delta_{\mathrm{tip}}$.

\textbf{Solution.}
\textbf{(a)} Center-to-center spacing of two $\SI{0.50}{\milli\meter}$ layers is $d = \SI{0.50}{\milli\meter} = \SI{5.00e-4}{\meter}$. From (\ref{eq:03:bimetallic}):
\[
  R_{\mathrm{curv}} \approx \frac{d}{(\alpha_2 - \alpha_1)\Delta T} = \frac{\SI{5.00e-4}{\meter}}{(19.0 - 1.2)\times 10^{-6}\,\si{\per\celsius} \times \SI{100.0}{\celsius}} = \frac{\SI{5.00e-4}{\meter}}{\num{1.780e-3}} = \mathbf{\SI{0.2809}{\meter} = \SI{28.1}{\centi\meter}}.
\]
\textbf{(b)} By the Pythagorean circle-chord relation $(R_{\mathrm{curv}} - \delta_{\mathrm{tip}})^2 + L_0^2 \approx R_{\mathrm{curv}}^2 \implies \delta_{\mathrm{tip}} \approx \dfrac{L_0^2}{2R_{\mathrm{curv}}}$ (zero trigonometry required!):
\[
  \delta_{\mathrm{tip}} \approx \frac{(\SI{0.0800}{\meter})^2}{2(\SI{0.28090}{\meter})} = \mathbf{\SI{0.01139}{\meter} = \SI{1.14}{\centi\meter}} \quad (\text{bending toward the Invar side}).
\]
\verifiedby{Exact formula (\ref{eq:03:bimetallic}) gives $R_{\mathrm{curv}} = \SI{0.28118}{\meter}$, $\delta_{\mathrm{tip}} = \SI{1.138}{\centi\meter}$ (\texttt{check\_sec02\_05.py}).}
\end{workedexample}

\begin{workedexample}[Joule's Falling-Weight Paddle-Wheel Experiment (\texttt{Rule 5.c.v})]{Competition $\cdot$ Rules 5.c.i, 5.c.v}{ex:03.5}
In a reproduction of James Prescott Joule's 1845--1850 paddle-wheel calorimeter experiment, two lead masses of $m = \SI{13.20}{\kilogram}$ each fall through a vertical height of $h = \SI{1.60}{\meter}$ a total of $N = 20$ times ($g = \SI{9.80665}{\meter\per\square\second}$), churning $\SI{5.50}{\kilogram}$ of water ($c_w = \SI{4184}{\joule\per\kilogram\per\kelvin}$) inside an insulated copper vessel of heat capacity $C_{\mathrm{cal}} = \SI{385}{\joule\per\kelvin}$. Assuming $100\%$ of the mechanical potential energy is dissipated into the water and vessel, calculate the expected temperature rise $\Delta T$ to $\SI{0.001}{\celsius}$ (explaining why Joule needed custom thermometers readable to $\frac{1}{200}\,^\circ\mathrm{F}$!).

\textbf{Solution.}
Total mechanical work delivered by the two falling masses over $N = 20$ drops is:
\[
  W_{\mathrm{mech}} = 2 N m g h = 2(20)(\SI{13.20}{\kilogram})(\SI{9.80665}{\meter\per\square\second})(\SI{1.60}{\meter}) = \mathbf{\SI{8284.66}{\joule}}.
\]
Total heat capacity of the water plus copper calorimeter is $C_{\mathrm{tot}} = m_w c_w + C_{\mathrm{cal}} = (5.50)(4184) + 385 = 23012 + 385 = \SI{23397}{\joule\per\kelvin}$. Thus:
\[
  \Delta T = \frac{W_{\mathrm{mech}}}{C_{\mathrm{tot}}} = \frac{\SI{8284.66}{\joule}}{\SI{23397}{\joule\per\kelvin}} = \mathbf{\SI{0.35409}{\kelvin} \approx \SI{0.354}{\celsius}} \quad (\text{or } 0.637\,^\circ\mathrm{F}).
\]
\verifiedby{$W_{\mathrm{mech}} / C_{\mathrm{tot}} = 8284.6579 / 23397.0 = \SI{0.354091}{\kelvin}$ (\texttt{check\_sec02\_05.py}).}
\end{workedexample}

\begin{trapbox}[§3 Top 5 Foundations, Thermometry \& Expansion Traps]
\begin{errtab}
Claiming a hole in a heated metal plate \emph{shrinks} as the metal expands inward & Sign of $\Delta D_{\mathrm{hole}}$ is inverted; a shrink-fit collar would jam instead of slipping on & Photographic enlargement rule: every linear distance (including cavities/holes) expands by $(1 + \alpha_L\Delta T)$. \\
Forgetting that the Delisle scale ($^\circ\mathrm{De}$) is \textbf{inverted} ($0\,^\circ\mathrm{De}$ boil, $150\,^\circ\mathrm{De}$ freeze) & Using $+1.5\,T_C$ instead of $150 - 1.5\,T_C$ gives the wrong sign and offset & Check fixed points: higher $^\circ\mathrm{De}$ means \textbf{colder} temperature ($\Delta T_{\mathrm{De}} = -1.5\,\Delta T_C$). \\
Confusing thermal expansion $\alpha_L$ ($\si{\per\kelvin}$) with thermal diffusivity $\alpha = k/(\rho c_p)$ ($\si{\square\meter\per\second}$) & Substitution of $\num{25e-6}$ where $\num{9.6e-5}$ belongs in transient conduction & Always check units: expansion is $\si{\per\kelvin}$; diffusivity is $\si{\square\meter\per\second}$. \\
Using $\beta_V = 3\alpha_L$ for liquid water between $\SI{0}{\celsius}$ and $\SI{4}{\celsius}$ & Predicts water expands upon warming from $\SI{0}{\celsius}$ to $\SI{4}{\celsius}$, when it actually \textbf{contracts} & Water has $\beta_V < 0$ for $\SI{0}{\celsius} < T < \SI{3.984}{\celsius}$ ($\rho_{\max} = \SI{999.97}{\kilogram\per\cubic\meter}$ at $\SI{3.984}{\celsius}$). \\
Attributing the 1780 ice calorimeter to Joseph Black or the $T^4$ theoretical proof to Stefan & Loses historical matching points under Rule \texttt{5.c.v} & Black discovered latent/specific heat (1761); Lavoisier \& Laplace built the ice calorimeter (1780). Stefan found $T^4$ empirically (1879); Boltzmann proved it theoretically (1884). \\
\end{errtab}
\end{trapbox}

\begin{devicebox}[Thermal Expansion \& Softening Limits Inside the 2027 Heat-Lamp Collector]
\idxdev{plastic party cup thermal softening and expansion!PET vs.\ PP}
Under a $\SIrange{100}{125}{\watt}$ heat lamp at $\SI{40}{\centi\meter}$ height, dry upper rim sections of the ES-provided 9~oz clear plastic party cup (\texttt{Rule 3.c, 4.g}) can reach $\SIrange{65}{80}{\celsius}$, while the water-wetted lower wall stays near $T_{\mathrm{water}} \approx \SIrange{30}{48}{\celsius}$. Clear party cups are made from \texttt{\#1 PET} ($T_g \approx \SI{74}{\celsius}$, $\alpha_L \approx \SI{65e-6}{\per\kelvin}$), \texttt{\#5 PP} ($\alpha_L \approx \SI{100e-6}{\per\kelvin}$), or \texttt{\#6 PS} (clear polystyrene, $\alpha_L \approx \SI{70e-6}{\per\kelvin}$, softens near $\SI{100}{\celsius}$). The ES's cup may be any of these: the polystyrene ban (FAQ \texttt{10/22/2026}) applies to \emph{your device}, not to the cup. Your cup cradle (\sref{sec:15}) must leave a $\ge\SI{1.5}{\milli\meter}$ radial clearance so thermal expansion never binds the cup during the 1-minute extraction window (\texttt{Rule 5.II.g}) \src{S01,S02}.
\end{devicebox}

\alsosee{\spref{sec:01.1} (\texttt{QS-01} Scales), \spref{sec:01.2} (\texttt{QS-02} Expansion), \spref{sec:01.14} (\texttt{QS-14} 8 Scientists), \spref{sec:04} (Calorimetry), \spref{sec:08} (Entropy \& 3rd Law), and \spref{sec:12} (Planck/Stefan--Boltzmann Radiation).}
